This is a small CPU study: tabular states, at most 6 rooms (155 cells), five seeds per cell and one hyperparameter setting shared by all systems. The tasks are deterministic apart from the TV token, the regime where optimistic initialisation alone explores systematically; the conclusion that the bonuses add nothing may not carry over to stochastic dynamics, function approximation or pixels, where optimism cannot be kept per state-action pair. The $\epsilon$-greedy baseline has no offset and is a random walk before the first reward; we did not tune $\epsilon$ or an optimistic $Q_0$ for it. The offset was introduced after a one-seed pilot. All systems are our own reimplementations. The RND-style bonus differs from the original in ways that matter here: the normaliser uses the standard deviation of raw errors instead of intrinsic returns, the bonus is combined with a constant offset, there is no separate intrinsic value head, and the guard ($W=64$, $c=5$) was added after the first version failed, although it was not tuned on results; predictor learning rate, update frequency and network size were not explored, and the clip sweep suggests a smaller clip would help a little. The $p$-values are uncorrected and use budget-censored values with $n=5$, and the penalty-only control, which cannot see the TV token, still differs between room\_4 and room\_4\_tv by more than its standard deviation, so they are descriptive. The noisy-TV test has one TV placement off the critical path, a finite uniform token set and no forward-model bonus, so it cannot reproduce the failure the noisy-TV argument describes. A full study would add stochastic dynamics, a forward-model bonus, a tuned optimistic baseline, more seeds and larger mazes.
